\clearpage
\phantomsection% 
\addcontentsline{toc}{chapter}{KATA PENGANTAR}
%\thispagestyle{fancy}

\begin{justifying}
	\large \bfseries \centering \MakeUppercase{Kata Pengantar}\linebreak
	
	\normalsize \normalfont \justifying
	Puji dan syukur penulis panjatkan ke hadirat Tuhan Yang Maha Esa atas berkat dan karunia-Nya, sehingga penyusunan laporan tugas akhir ini dapat diselesaikan dengan baik. Dalam proses penyusunan laporan tugas akhir ini, penulis menyadari bahwa keberhasilan penyelesaian penelitian tidak terlepas dari bantuan, bimbingan, arahan, serta dukungan dari berbagai pihak. Oleh karena itu, pada kesempatan ini penulis ingin menyampaikan ucapan terima kasih yang sebesar-besarnya kepada: \par

	\begin{enumerate}
		\item Rasa syukur yang tak terhingga penulis haturkan kepada orang tua dan keluarga tercinta, yang senantiasa memberikan doa, kasih sayang, dukungan moral, serta motivasi tanpa henti selama penulis menempuh pendidikan hingga penyusunan tugas akhir ini dapat diselesaikan.
		\item Bapak I Wayan Wiprayoga Wisesa, S.Kom., M.Kom. selaku Dosen Pembimbing, atas bimbingan, arahan, perhatian, serta masukan yang sangat berharga selama proses penelitian dan penyusunan tugas akhir ini.
		\item Bapak Hafiz Budi Firmansyah S.Kom., M.Sc., Ph.D. selaku Dosen Penguji I dan Bapak Rahman Indra Kesuma, S.Kom., M.Cs. selaku Dosen Penguji II, yang telah memberikan saran, kritik, serta masukan yang konstruktif demi penyempurnaan tugas akhir ini.
		\item Bapak Radhinka Bagaskara, S.Si.Kom., M.Si., M.Sc. selaku Dosen Wali, atas bimbingan, arahan, serta dukungan yang diberikan kepada penulis selama menempuh masa perkuliahan di Program Studi Teknik Informatika.
		\item Seluruh dosen dan staf pengajar Program Studi Teknik Informatika, Institut Teknologi Sumatera, yang telah memberikan ilmu pengetahuan, pengalaman, serta wawasan yang sangat bermanfaat bagi penulis selama masa studi.
		\item Teman-teman Gazebo (Lathief, Apridian, Alfajar, Joshua, Bezalel, Ridho, Irma, Keti, dan Shintya), yang telah menjadi bagian dari perjalanan akademik penulis, memberikan dukungan, kebersamaan, dan semangat selama masa perkuliahan.
		\item Seluruh pihak lain yang tidak dapat disebutkan satu per satu, yang telah memberikan bantuan, baik secara langsung maupun tidak langsung, dalam proses penyusunan tugas akhir ini.
	\end{enumerate} \par

	Penulis menyadari bahwa laporan tugas akhir ini masih memiliki keterbatasan dan jauh dari kata sempurna. Oleh karena itu, penulis dengan terbuka menerima kritik dan saran yang membangun sebagai bahan evaluasi dan perbaikan di masa mendatang. \par

	Akhir kata, penulis berharap semoga laporan tugas akhir ini dapat memberikan manfaat dan kontribusi positif bagi pengembangan ilmu pengetahuan dan teknologi, khususnya di bidang Teknik Informatika, serta dapat menjadi referensi untuk penelitian selanjutnya.
	
\end{justifying}
\clearpage
